% from ega4-errata-addenda-list3-fr.tex
\oldpage[IV]{350}
\paragraph{(Err\textsubscript{IV}, 20).}
Après (IV, 4.5.19), insérer, avant la ligne 12 du bas de la p.~67 :

\begin{corollary}[4.5.19.3]
\label{IV.4.5.19.3}
Soient $k$ un corps, $X$ un $k$-préschéma de type fini sur $k$, qui est
irréductible mais non géométriquement irréductible. Alors il existe une
partie fermée $Y$ de $X$, rare dans $X$, contenant tous les points
$x\in X$ tels que $k(x)$ soit radiciel sur $k$ (et en particulier tous
les points de $X$ rationnels sur $k$).
\end{corollary}

Soient $k'$ une clôture algébrique de $k$, $X'=X\otimes_kk'$,
$p:X'\to X$ la projection canonique. Si $x\in X$ est tel que $k(x)$
soit radiciel sur $k$, il n'existe qu'un seul point $x'\in X'$ au-dessus
de $x$ (I, 3.4.9); ce point appartient nécessairement à deux composantes
irréductibles distinctes de $X'$, sans quoi $X$ serait géométriquement
irréductible en vertu de
{\hyperref[canon-index.ega4.c04.prop.4-5-19]{(4.5.19)}}, contrairement à
l'hypothèse. Soient alors $X'_i$ ($1\leq i\leq n$) les composantes
irréductibles du préschéma noethérien $X'$, et pour
$1\leq i<j\leq n$, posons $X'_{ij}=X'_i\cap X'_j$; l'ensemble
\[
 Y=p\left(\bigcup_{i<j}X'_{ij}\right)
\]
répond à la question. En effet, comme $p$ est un morphisme entier et
que les $X'_{ij}$ sont fermés dans $X'$, $Y$ est fermé dans $X$.
D'autre part, si $\xi$ est le point générique de $X$, $p^{-1}(\xi)$
est discret et formé des points génériques des $X'_i$
{\hyperref[canon-index.ega4.c04.cor.4-5-11]{(4.5.11)}}, donc ne rencontre aucun
des $X'_{ij}$; l'ensemble fermé $Y$ est par suite rare.

\paragraph{(Err\textsubscript{IV}, 21).}
Dans (IV, 5.6.11.1), ligne 7 du bas de la p.~102, après « intègre »,
ajouter : « de dimension 2 ».

\paragraph{(Err\textsubscript{IV}, 22).}
Après (IV, 6.2.3), après la ligne 18 de la p.~138, insérer :

\begin{proposition}[6.2.4]
\label{IV.6.2.4}
Sous les hypothèses de {\hyperref[canon-index.ega4.c04.prop.6-2-3]{(6.2.3)}}, soit
de plus $M$ un $A$-module de type fini et supposons $M\neq0$ et
$N\neq0$. Alors on a
\begin{equation}
 \operatorname{dim.proj}_B(M\otimes_AN)=
 \operatorname{dim.proj}_A(M)+
 \operatorname{dim.proj}_{B\otimes_Ak}(N\otimes_Ak).
 \label{IV.6.2.4.1}\tag{6.2.4.1}
\end{equation}
\end{proposition}

Considérons la résolution gauche $L_\bullet$ de $N$ introduite dans
{\hyperref[canon-index.ega4.c04.prop.6-2-3]{(6.2.3)}}, et une résolution gauche
$K_\bullet$ de $M$ par des $A$-modules libres de type fini :
\[
 \cdots\longrightarrow K_i\longrightarrow K_{i-1}\longrightarrow
 \cdots\longrightarrow K_0\longrightarrow M\longrightarrow0.
\]
On a vu que les $L_i$, $N$, les $Z_i(L_\bullet)$ sont des $A$-modules
plats; il en est de même des $B_i(L_\bullet)$ qui sont égaux aux
$Z_i(L_\bullet)$ sauf pour $i=0$, et dans ce dernier cas
$B_0(L_\bullet)$ est

